\subsection{\texorpdfstring{\(\mathfrak{S}\)-收敛的例子}{S-收敛的例子}}

I. 在 X 的子集上的一致收敛。设 A 是 X 的子集，取 \(\mathfrak{S} = \{\mathrm{A}\}\)。\(\mathfrak{S}\)-收敛的一致结构（相应地，拓扑）这时称为在 A 中一致收敛的一致结构（相应地，拓扑）；若滤子 \(\Phi\)（\(\mathcal{F}_{\mathfrak{S}}(\mathrm{X};\mathrm{Y})\) 上的）收敛于 \(u_0\)，则称它在 A 中一致收敛于 \(u_0\)。当 \(\mathrm{A} = \mathrm{X}\) 时，我们重新得到第 1 号中定义的一致收敛结构。

\begin{enumerate}
\def\labelenumi{\Roman{enumi}.}
\setcounter{enumi}{1}
\tightlist
\item
  在 X 的子集上的逐点收敛。设 A 是 X 的子集，取 S 为由属于 A 的单个点组成的 X 的所有子集之集（由第 2 号的注记 2，若取 S 为 A 的所有有限子集之集，效果相同）。S-收敛的一致结构（相应地，拓扑）这时称为在 A 中逐点收敛的一致结构（相应地，拓扑）；若 \(\mathcal{F}_{\mathfrak{S}}(\mathbf{X};\mathbf{Y})\) 上的滤子 Φ 收敛于 \(u_0\)，则称它在 A 中逐点收敛于 \(u_0\)。这种情形成立当且仅当：对每个 \(x\in \mathbf{A}\)，\(u_0(x)\) 是 \(u(x)\) 关于滤子 Φ 的极限。
\end{enumerate}

特别地，当 A = X 时，在 X 中逐点收敛的一致结构（相应地，拓扑）就简称为逐点收敛的一致结构（相应地，拓扑）；把 \(\mathcal{F}(X; Y)\) 赋以这结构所得的一致空间记作 \(\mathcal{F}_{s}(X; Y)\)。注意逐点收敛的拓扑就是 \(Y^{X}\) 上的乘积拓扑，因而只依赖于 Y 的拓扑而不依赖其一致结构。

\begin{enumerate}
\def\labelenumi{\Roman{enumi}.}
\setcounter{enumi}{2}
\tightlist
\item
  紧收敛。设 X 是拓扑空间，取 S 为 X 的所有紧子集组成的集。这时 S-收敛的一致结构（相应地，拓扑）称为紧收敛的一致结构（相应地，拓扑），而赋予 \(\mathcal{F}(X;Y)\) 此一致结构所得的一致空间记作 \(\mathcal{F}_{c}(X;Y)\)。紧收敛的结构比一致收敛的结构粗，且当 X 紧时二者一致；它又比逐点收敛的结构细，且当 X 离散时这二者一致。
\end{enumerate}

若 X 是一致空间，则通过取 S 为 X 的所有准紧子集组成的集，我们可以在 \(\mathcal{F}(X;Y)\) 上定义准紧收敛的一致结构。同样，若 X 是度量空间，则可取 S 为 X 的所有有界子集组成的集；这时 S-收敛的一致结构称为有界收敛的一致结构。

